% Generated draft: requires mathematical and source review.
\documentclass[11pt]{article}
\usepackage[margin=25mm]{geometry}
\usepackage{fontspec}
\setmainfont{Latin Modern Roman}
\usepackage{amsmath,amssymb,mathtools,unicode-math}
\setmathfont{Latin Modern Math}
\usepackage{enumitem,needspace,xurl,hyperref}
\setlength{\parindent}{0pt}
\setlength{\parskip}{5pt}
\setlength{\emergencystretch}{3em}
\newcommand{\N}{\mathbb{N}}
\newcommand{\Z}{\mathbb{Z}}
\newcommand{\Q}{\mathbb{Q}}
\newcommand{\R}{\mathbb{R}}
\newcommand{\C}{\mathbb{C}}
\newcommand{\F}{\mathbb{F}}
\newcommand{\A}{\mathbb{A}}
\newcommand{\PP}{\mathbb{P}}
\newcommand{\E}{\mathbb{E}}
\newcommand{\eps}{\varepsilon}
\newcommand{\dd}{\,\mathrm{d}}
\begin{document}
\section*{UnsolvedMath ID: 30006016}
\textbf{Source problem code:} catalogue entry
\subsection{Definitions and mathematical statement}
\textbf{Supplied research problem.} Extend the existing bijective approach to self-overlapping polygons beyond its present topological setting to other surface topologies, and characterize the enumerative formulas arising from such extensions.

A precise extension would have to specify the original class of self-overlapping polygons, the objects in bijection with that class, the original surface topology, the additional surface topologies under consideration, and the corresponding bijective constructions. Any enumerative formulas would also require specified counting parameters, equivalence relations, and weights or marking conventions, where applicable. None of these data is defined in the supplied record. In particular, the record does not specify genus, orientability, boundary conditions, or whether polygons are considered in a fixed surface. Thus the request does not supply a quantified assertion that an extension exists for every surface topology.

\textbf{Unresolved source qualifications.} The quoted source sentence asks whether the authors' bijective approach can be extended to other topologies. The accompanying extracted formula is corrupted: it contains notation resembling \(W_n^{(g)}\), \(T_{\ell_1,\ldots,\ell_n}^{(g)}\), \(x(z_i)\), \(x'(z_i)\), \(\gamma\), derivatives, and indices \(\ell_1,\ldots,\ell_n\geq 1\), but gives neither reliable formula structure nor definitions of these symbols. Consequently, no identity or counting interpretation is asserted here. The additional request to characterize enumerative formulas appears in the supplied cleaned statement; its precise connection to the quoted sentence and damaged formula requires source review. This is a statement-only draft of an incompletely specified research question.
\subsection{Short English statement}
The problem asks whether an existing bijection for self-overlapping polygons extends to other surface topologies and what counting formulas such extensions yield. The supplied record leaves the original bijection, the admissible polygons, the surface conditions, and the counting conventions unspecified, so a precise theorem statement requires further source review.
\subsection{Sources}
\begin{itemize}
\item Local UnsolvedMath registry, \texttt{data/unsolvedmath/problems.json}, permanent ID \(30006016\), source problem code \texttt{catalogue entry}, title \emph{Bijections for Self-Overlapping Polygons in Other Topologies}. The supplied record belongs to the Oberwolfach Reports Open Problems collection. No individual proposer is identified; the original quotation refers to an unattributed authors' approach.
\item Supplied report citation: \emph{Riemann Surfaces: Random, Flat, and Hyperbolic Geometry} (2025), Oberwolfach Reports, pp.~2369--2424, DOI \texttt{10.4171/owr/2024/41}. \url{https://doi.org/10.4171/owr/2024/41}. The record describes Oberwolfach Reports as documenting mathematical workshops and their open-problem discussions. It records proposed year 2024, whereas the supplied report citation gives 2025; both are retained without resolving their bibliographic relationship.
\item Supplied literature source labelled Shor--van Wyk: \url{https://epubs.siam.org/doi/10.1137/1.9781611973068.19}. No title, theorem, or explanation of its exact relationship to the proposed extension is supplied.
\item Literature assessment dated exactly \texttt{2026-08-22}: the supplied curation labels the problem open, states that the report presents extension beyond the stated topology as an open question, and reports that no later primary result located in that curation resolves the general extension. Its listed sources checked are the OWR report and Shor--van Wyk reference above. This is supplied dated literature triage derived from OWR curation, not an independently verified assessment; time-sensitive status and cited claims require independent verification.
\end{itemize}
\end{document}
